\documentclass[11pt,a4paper]{article}
\usepackage[T1]{fontenc}
\usepackage{isabelle,isabellesym}

\usepackage{amsmath}
\usepackage{amssymb}
\usepackage{stmaryrd}

% this should be the last package used
\usepackage{pdfsetup}

\urlstyle{rm}
\isabellestyle{it}

\begin{document}

\title{Sup-Convolutions, Alexandrov's Theorem, and Comparison of Viscosity Solutions}
\author{Dmitriy Traytel}
\maketitle

\begin{abstract}
  Comparison between a viscosity subsolution and a supersolution of a fully
  nonlinear degenerate elliptic equation rests on second-order
  differentiability of functions that are merely semiconvex.  This entry
  brings that theory to Isabelle.  We formalize Rademacher's theorem, that a
  locally Lipschitz function on a Euclidean space is differentiable at almost
  every point, and Alexandrov's theorem, that a convex function has a
  second-order Taylor expansion at almost every point, over a base of
  subgradients, the proximal map, the Moreau envelope and the
  sup-convolution; on them rest Jensen's lemma for semiconvex functions and
  the Crandall--Ishii theorem on sums, which at a local maximum of
  $u(x)+w(y)$ produces second-order jets for $u$ and $w$ whose Hessians
  satisfy a matrix inequality.  Two further theories supply the
  doubling-of-variables apparatus that consumes these results, together with
  coercive penalties whose second-order jet vanishes on the diagonal.  The
  theorem on sums is also stated in the vocabulary of the literature:
  second-order semijets and their closures, with two named theorems for the
  quadratic coupling $\frac{\alpha}{2}\lVert x-y\rVert^2$.  The staged form
  exhibits sequences of genuine ordered jets with quantitative control; the
  closed-jet form places, for every scale $\lambda$ with $2\lambda\alpha<1$,
  a pair $X \preceq Y$ with
  $-\frac{1}{\lambda}I \preceq X \preceq Y \preceq \frac{1}{\lambda}I$ in
  the closed semijets at the maximum point itself, the shift of the
  maximiser under the deconvolved penalty
  $\frac{\alpha}{1-2\lambda\alpha}$ cancelling the sup-convolution's
  displacement of the jets exactly.
  Alexandrov's theorem is obtained by Minty's device, applying Rademacher to
  the resolvent of the subdifferential and transporting the expansion back
  from the Moreau envelope; the doubling toolbox is stated over an arbitrary
  Euclidean space, matrices appearing only where a statement produces one.
  We follow Evans and Gariepy~\cite{EvansGariepy} and the User's Guide of
  Crandall, Ishii and Lions~\cite{CrandallIshiiLions}.
\end{abstract}

\tableofcontents

\parindent 0pt\parskip 0.5ex

\input{session}

\bibliographystyle{abbrv}
\bibliography{root}

\end{document}
